\documentclass{article}
\usepackage{polski}
\usepackage[a4paper, margin=2.5cm]{geometry}
\usepackage{graphics}
\usepackage[export]{adjustbox}
\usepackage{float}
\usepackage{amsmath}
\usepackage{xcolor}
\usepackage{multirow}
\usepackage{array}


\title{Analiza wyników pomiarowych z prób spęczania z wykorzystaniem aplikacji komputerowej}
\date{12.01.2026}
\author{Szymon Miłkowski}

\begin{document}
	\maketitle
	
	\section{Cel i wstępne założenia}
	Celem tego projektu było stworzenie graficznego programu komputerowego do analizy, korekcji i transformacji danych pomiarowych z prób spęczania. Danymi wejściowymi do programu miały być pliki binarne w formacie \textit{W01} generowane przez oprogramowanie maszyny badawczej. \\
	Dodatkowym celem miała być kompatybilność i płynność działania zarówno na współczesnym sprzęcie, jak i starszym (Windows XP 32-bit).

	\section{Wybór technologii}
	Z racji na wymóg szerokiej kompatybilności, oraz wysokiej wydajności, konieczne było wybranie technologii pozwalającej na kompilację do natywnego kodu maszynowego z minimalną ilością zależności.
	Wybór padł na język C z biblioteką Raylib do obsługi grafiki, linkowaną statycznie.\\
	Raylib posiada łatwe w obsłudze API składające się z funkcji do rysowania i obsługi wejścia (mysz, klawiatura).
	
	\section{Spęczanie}
	Próba spęczania polega na ściskaniu próbki wzdłuż jednej osi (pionowej) przez maszynę badawczą z zadaną prędkością. Podczas próby materiał odkształca się, zmniejszając swoją wysokość w kierunki siły oraz - zgodnie z zasadą zachowania objętości - zwiększając swoją szerokość w pozostałych kierunkach. \\
	Podszas tego procesu, materiał wewnątrz płynie, wypychając materiał na ścianach bocznych na zewnątrz, a materiał przy podstawach trze o górną i dolną ścianę maszyny, co hamuje jego przemieszczenie na zewnątrz. \\
	Siła nacisku (opór jaki stawia próbka) jest mierzona i zależy od materiału i fazy procesu.
	
	\subsection{Fazy procesu}
	Brak kontaktu
	\begin{itemize}
		\item Brak kontaktu - maszyna nie ściska próbki, zmierza w jej kierunku - stała, niska siła
		\item Pierwszy kontakt - luzy w maszynie są ściskane - szum informacyjny
		\item Faza liniowa?? - siła rośnie liniowo
		\item Faza nieliniowa?? - siła rośnie coraz mniej
	\end{itemize}
	
	
	
	\pagebreak
	
	
	\section{Format pliku wejściowego}
	\subsection{Nieznany format binarny}
	Pliki wejściowe z danymi pomiarowymi z rozszerzeniem \textit{.W01}, zapisane są w nieudokumentowanym formacie binarnym. Wpisywanie w wyszukiwarce \textit{.W01} lub magic number z pliku - pierwsze cztery bajty \textit{66 66 86 3f} nie daje istotnych rezultatów. Aby zczytać dane z plików i stworzyć algorytm importujący do programu, należy więc przeprowadzić inżynierię wsteczną formatu.
	
	\subsection{Inżynieria wsteczna formatu}
	W inżynierii wstecznej nie ma jednego poprawnego przepisu na znalezienie rozwiązania. Każde zadanie jest osobną łamigłówką, do której musimy użyć narzędzi, doświadczenia, wiedzy informatycznej, wiedzy dziedzinowej oraz czasu.
	Podczas analizy od dyspozycji miałem dwa pliki: \textit{19\_.v10.W01} oraz \textit{20\_.v50.W01}, gdzie \textit{v10} i \textit{v50} odpowiadają prędkości procesu spęczania.
	
	\subsubsection{Pierwsze spojrzenie na plik}
	Rzut okiem w \textit{hex-edytorze} ujawnia kilka stringów i sporo zer na początku pliku, a w dalszej części powtarzający się czterokrotnie wzór: napis i blok nieczytelnych danych.
	
	\subsubsection{Nagłówek}
	W nagłówgu pliku, który znajduje się na początku, przed właściwymi danymi liczbowymi znajdują się metadane, informacje o rozmiarze, godzina utworzenia oraz nazwa programu wejściowego do maszyny badawczej. Cały nagłówek ma stały rozmiar 208 bajtów.
	
	\begin{center}
		\begin{tabular}{|c|c|c|c|p{6.5cm}|} \hline
			\multicolumn{5}{|c|}{Struktura nagłówku pliku} \\\hline
			offset & rozmiar & wartość/typ & nazwa & komentarz \\\hline
			\hline
			 0 &  4 & 0x3f866666 & magic & stała wartość określająca format \\\hline
			 4 &  4 & 1 & wersja? & najpewniej wersja formatu \\\hline
			 8 &  4 & $rozmiar\_pliku + 1$ & rozmiar & \\\hline
			12 &  4 & 0 & nieznana & zawsze 0 \\\hline
			16 &  8 & "ScriptPl" & stały napis & 8 znaków, bez null terminatora \\\hline
			24 &  4 & $rozmiar\_pliku - 16$ & rozmiar2 & wygląda jak zbędna informacja, może oznaczać "rozmiar począwszy od napisu" \\\hline
			28 &  4 & liczba wierszy & liczba wierszy & ilość liczb w każdej kolumnie danych \\\hline
			32 &  4 & liczba kolumn & liczba kolumn & ilość kolumn z danymi \\\hline
			36 &  4 & 1 & nieznane & \\\hline
			40 &  4 & 1 & nieznane & \\\hline
			44 & 18 & "dd-MM-yyyyhh:mm:ss" & czas utworzenia & napis bez null terminatora \\\hline
			62 & 42 & nieznana & nieznane & 42 bajty, głównie zera \\\hline
			104 & 104 & nazwa pliku & plik badania & nazwa pliku wejściowego do maszyny, dopełniona zerami do 104 bajtów \\\hline
		\end{tabular}
	\end{center}
	
	\subsubsection{Kolumny (serie danych)}
	Kolumny mają zmienny rozmiar, zależnie od pliku, jednak zawsze taki sam w obrębie jednego pliku.
	
	\begin{center}
		\begin{tabular}{|c|c|c|c|c|} \hline
			\multicolumn{5}{|c|}{Struktura kolumny} \\\hline
			offset & rozmiar & wartość/typ & nazwa & komentarz \\\hline
			\hline
			0 & 4 & int & nieznane & 1 lub 16385  \\\hline
			4 & 59 & napis & nazwa & nazwa i jednostka kolumny, oddzielone spacjami \\\hline
			63 & $4 * liczba\_wierszy$ & float & dane & tablica 32-bitowych liczb zmiennoprzecinkowych \\\hline
		\end{tabular}
	\end{center}
	
	\pagebreak
	\subsubsection{Struktura całego pliku}
	Plik składa się z nagłówka o stałym rozmiarze i zmiennej liczby kolumn zmiennej długości, które zawierają swój nagłówek oraz danych liczbowych.
	
	\begin{center}
		\begin{tabular}{|p{5cm}|p{5cm}|} \hline
			\multicolumn{2}{|c|}{Struktura pliku} \\\hline
			rozmiar  & element \\\hline
			\hline
			208 & nagłówek  \\\hline
			$63+4*liczba\_wierszy$  & kolumna 1 \\\hline
			$63+4*liczba\_wierszy$  & kolumna 2 \\\hline
			...  & ... \\\hline
			$63+4*liczba\_wierszy$  & kolumna $n$ \\\hline
		\end{tabular}
	\end{center}
	
	
	\pagebreak
	
	

	\section{Ogólna struktura aplikacji}
	\subsection{Komponenty}
	Interfejs graficzny aplikacji podzielony jest na "komponenty" - funkcje zawierające logikę kontrolną oraz rysujące pewien element interfejsu. Wywołanie tej funkcji oznacza, że komponent istnieje w obecnej klatce.
	Funkcje te mogą wywoływać inne funkcje komponentów, co tworzy swego rodzaju strukturę drzewa komponentów, rozpoczynającą się w głównej pętli programu. Pozycja w strukturze drzewa nie ma jednak wpływu na rozmieszczenie elementu, pozycja jest statyczna, lub określana poprzez podane parametry.
	
	\subsection{Zarządzanie stanem}
	Inspirowane działaniem innego drzewiastego systemu komponentów graficznych - webowym Reactem - komponenty, jeśli to konieczne, operują na dwóch rodzajach danych: właściwości (\textit{props}) i stan (\textit{state}).
	\paragraph{Właściwości} są podawane do funkcji komponentu jako kopia (pass by value), bądź stały wskaźnik, mają za zadanie sterować zachowaniem i wyglądem komponentu.
	\paragraph{Stan} jest podawany jako wskaźnik do pamięci zarządzanej przez rodzica (komponent wywołujuący), ma za zadanie przechowywanie stanu komponentu między wywołaniami, jest modyfikowany głównie przez sam komponent oraz jego dzieci (jeśli zostanie podany do nich jako ich stan).
	\paragraph{Stan globalny} jest wykorzystywany w niektórych komponentach, które z założenia mają istnieć tylko w jednej sztuce jednocześnie, jak lista otwartych plików lub narzędzia diagnostyczne \textit{FPSCounter} i \textit{Clock}, jako zmienne globalne w jednostce translacyjnej (niewidoczno poza obecnym plikiem .c), lub statyczne zmienne lokalne w funkcji komponentu.
	
	
	
	\section{Rysowanie wykresu}
	Raylib zawiera jedynie proste funkcje rysowania kształtów. Jedną z nich jest \textit{DrawSplineLinear}, łącząca podane współrzędne pikseli liniowymi odcinkami. cośtam właśnie jej użyjemy bla
	
	\begin{verbatim}
		DrawSplineLinear(state->point_buffer, state->visible.count, 2.f, DARKBLUE);
	\end{verbatim}
	
	\subsection{Poruszanie się w widoku}
	Aby zwiększyć interaktywność, użytkownik może się poruszać po płaszczyźnie XY, przeciągając płaszczyznę myszką, oraz przybliżając i oddalając kółkiem myszy. Manipulowane są w ten sposób zmienne skali oraz przesunięcia uwzględnione w przekształceniu punktów ze współrzędnych kartezjańskich na współrzędne ekranu.
	
	\subsection{Transformacja współrzędnych}
	
	\subsubsection{Model matematyczny}
	Punkt podlega przekształceniu liniowemu:

	\begin{equation}
		x = \frac{z(x_p - x_v) - x_{S}}{x_{E} - x_{S}} * W + x_O
	\end{equation}
	gdzie \\
	$z$ - współczynnik skali \\
	$x_p$ - współrzędna $x$ punktu w układzie kartezjańskim \\
	$x_v$ - przesunięcie $x$ w układzie kartezjańskim \\
	$x_S$ - współrzędna $x$ lewej krawędzi układu kartezjańskiego przy zerowym przesunięciu \\
	$x_E$ - współrzędna $x$ prawej krawędzi układu kartezjańskiego przy zerowym przesunięciu \\
	$W$ - szerokość wykresu w pikselach \\
	$x_O$ - odległość prawej krawędzi wykresu od prawej krawędzi ekranu
	\\\\
	
	\begin{equation}
		y = \frac{-z(y_p + y_v) - y_{S}}{y_{E} - y_{S}} * H + y_O
	\end{equation}
	gdzie \\
	$z$ - współczynnik skali \\
	$y_p$ - współrzędna $y$ punktu w układzie kartezjańskim \\
	$y_v$ - przesunięcie $y$ w układzie kartezjańskim \\
	$y_S$ - współrzędna $y$ dolnej krawędzi układu kartezjańskiego przy zerowym przesunięciu \\
	$y_E$ - współrzędna $y$ górnej krawędzi układu kartezjańskiego przy zerowym przesunięciu \\
	$H$ - wysokość wykresu w pikselach \\
	$y_O$ - odległość górnej krawędzi wykresu od górnej krawędzi ekranu \\
	
	Zmiana znaku $y_p$ jest spowodowana, tym że oś +Y we współrzędnych ekranowych skierowana jest w dół.
	
	\subsubsection{Implementacja}
	W kodzie programu zaimplementowane jest to dwuetapowo.
	\paragraph{Obliczanie granic widocznego fragmentu}
	\begin{verbatim}
		Bounds compute_bounds(const float zoom, const Vector2 pan) {
			return (Bounds) {
				.start_x = PLOT_START_X / zoom + pan.x,
				.end_x = PLOT_END_X / zoom + pan.x,
				.start_y = PLOT_START_Y / zoom / ((float) PLOT_WIDTH / PLOT_HEIGHT) + pan.y,
				.end_y = PLOT_END_Y / zoom / ((float) PLOT_WIDTH / PLOT_HEIGHT) + pan.y,
			};
		}
	\end{verbatim}
	
	\paragraph{Transformacja pojedynczego punktu}
	\begin{verbatim}
		Vector2 transform_v_to_pixel(const Vector2 v, const Bounds bounds) {
			return (Vector2){
				.x = (v.x - bounds.start_x) / (bounds.end_x - bounds.start_x) * PLOT_WIDTH + PLOT_OFFSET_X,
				.y = (-v.y - bounds.start_y) / (bounds.end_y - bounds.start_y) * PLOT_HEIGHT + PLOT_OFFSET_Y,
			};
		}
	\end{verbatim}
	
	\paragraph{Transformacja wielu punktów jednocześnie} jest zoptymalizowana pod kątem ilości niepotrzebnych operacji (opisana w punkcie 6.3.3).
	
	\subsection{Minimalizacja renderowanych punktów}
	\subsubsection{Opis problemu}
	Podczas przegladania wykresu, nie każdy punkt będzie widoczny jednocześnie. Chcąc zminimalizować ilość punktów podlegających transformacji oraz rysowaniu, należy znaleźć zakres punktów widocznych na ekranie.
	
	\subsubsection{Szukanie pierwszego punktu}
	Zakładając, że współrzędna $x$ w zestawie danych jest niemalejąca, możemy potraktować zbiór $X$ jako posortowany zbiór danych i użyć wydajnego algorytmu \textbf{wyszukiwania binarnego}, o złożoności obliczeniowej $O(\log_2n)$.
	
	\begin{verbatim}
		int find_first_visible(const Vector2* data, const int count, const float start_x) {
			if (data[count - 1].x < start_x) return -1;
			
			int a = 0, b = count - 1;
			while (b - a > 1) {
				const int mid = (a + b) / 2;
				if (data[mid].x > start_x) b = mid;
				else a = mid;
			}
			return a;
		}
	\end{verbatim}
	
	\subsubsection{Transformacja punktów}
	Po znalezieniu pierwszego elementu zbioru wyjściowego, współrzędne są przekształcane i zapisywane w tablicy wyjściwej aż do napotkania punktu znajdującego się za prawą krawędzią widoku lub końca tablicy wejściowej. Wyszukiwanie binarne nie ma tutaj zastosowania, ponieważ i tak każdy widoczny punkt musi zostać przekształcony, co wiąże się z liniowym algorytmem.
	
	\begin{verbatim}
		VisiblePointsInfo compute_visible_points_offset(
			const DataSource data_source, Vector2* point_buffer, const Bounds bounds, const Vector2 offset
			) {
			VisiblePointsInfo res = {
				.start = find_first_visible(data_source.data, data_source.count, bounds.start_x - offset.x)
			};
			if (res.start == -1) return res;
			
			const float x_factor =  1 / (bounds.end_x - bounds.start_x) * PLOT_WIDTH;
			const float y_factor = -1 / (bounds.end_y - bounds.start_y) * PLOT_HEIGHT;
			
			const int visible_count = data_source.count - res.start;
			for (int i = 0; i < visible_count; i++) {
				const int i_source = i + res.start;
				point_buffer[i].x = (data_source.data[i_source].x - bounds.start_x + offset.x)
					* x_factor + PLOT_OFFSET_X;
				point_buffer[i].y = (data_source.data[i_source].y + bounds.start_y - offset.y)
					* y_factor + PLOT_OFFSET_Y;
				if (data_source.data[i_source].x > bounds.end_x - offset.x) {
					res.count = i + 1;
					return res;
				}
			}
			res.count = visible_count;
			return res;
		}
	\end{verbatim}
	Algorytm nie używa funkcji do transformacji pojedynczego punktu, co pozwala na obliczenie niezmiennych współczynników poza pętlą, unikając zbędnych operacji.
	
	\subsection{Wykrywanie zmian}
	Transformacja wszystkich widocznych punktów jest mimo wszystko kosztowną operacją, dlatego wykonywana jest jedynie gdy reprezentacja punktów na wykresie powinna ulec zmianie. Odpowiada za to zmienna \textit{change}, zawierająca flagi sterujące rekalkulacją elementów wykresu, ustawiana przez funkcję odpowiadającą za obsługę wejścia.
	
	
	\pagebreak
	\section{Analiza danych spęczania}
	\subsection{Charakterystyka danych}
	Dane zawarte w plikach z pomiaru próby spęczania zawierają 4 kolumny
	\subsubsection{Load[kN]}
	Siła zmierzona przez maszynę - opór, który stawia próbka. Wartości ujemne.
	\subsubsection{Stroke[mm]}
	Przemieszczenie narzędzia maszyny. Wartości ujemne.
	\subsubsection{Command[\%]}
	
	\subsubsection{Time[s]}
	Czas obecnego rekordu, wartość rosnąca od 0, co 0.0025.
	
	
	\subsection{Fazy procesu w danych}
	\begin{itemize}
		\item \textbf{Brak kontaktu} - obciążenie jest stałe lub z minimalnymi wachaniami, niewielkie
		\item \textbf{Pierwszy kontakt} - obciążenie wacha się nierównomiernie, zauważalnie bardziej niż w poprzednim etapie
		\item \textbf{Liniowa} - obciążenie rośnie liniowo
		\item \textbf{Nieliniowa} - wzrost obciążenia stopniowo i płynnie maleje
	\end{itemize}
	
	\subsection{Wykrywanie faz}
	Fazy procesu są na tyle charakterystyczne, że można je wykryć z poziomu algorytmu
	\\do napisania jeszcze
	
	\subsection{Analiza fazy liniowej}
	
	\subsection{Analiza fazy nieliniowej}
	
	
	\pagebreak
	
	
	\section{Proces optymalizacji algorytmu regresji wielomianowej}
	\subsection{Wstęp}
	Ten rozdział opisuje proces opytmalizacji algorytmu wykonującego regresję wielomianową podanego stopnia dla podanych punktów. Udokumentowane są użyte techniki oraz pomiary czasowe. Przed opytmalizacją, algorytm ten był zdecydowanie najbardziej czasochłonną operacją w całym programie.
	\subsection{Dane wejściowe, sprzęt i środowisko}
	Podczas testowania przyjęto następujące parametry:
	\begin{itemize}
		\item \textbf{Punkty wejściowe}: 29 punktów na krzywej $2\sin x+2$, dla $x$ od -5, co 0.3
		\item \textbf{Stopień szukanego wielomianu}: 28
		\item \textbf{CPU}: AMD Zen2 Ryzen 3600
		\item \textbf{Kompilator}: GCC 15.2.1
		\item \textbf{Opcje kompilatora w trybie Debug}: -g
		\item \textbf{Opcje kompilatora w trybie Release}: -O3
	\end{itemize}
	
	\subsection{Sposób pomiarów}
	\subsubsection{Realny benchmark}
	Pomiary algorytmu oraz jego fragmentów miały miejsce jako mierzenie części całego działającego programu, maksymalnie raz na klatkę, przy zachowaniu całej jego pozostałej funkcjonalności. Rezultat był wykorzystany do wyświetlenia na żywo wykresu wyznaczonej funkcji, nie ma mowy więc o pomijaniu przez kompilator operacji, których wynik nie jest używany. Dodatkowo w syntetycznym benchmarku, wywołującym w pętli funkcję wiele razy, korzystniej mogłaby być zapełniona pamięć cache oraz bardziej dokładny predyktor skoków, dlatego wplecenie pomiarów w standardowy przepływ programu lepiej oddaje wydajność w realnych warunkach.
	
	\subsubsection{RDTSC}
	Architektura x86 implementuje instrukcję \textbf{RDTSC} (\textit{Read Time-Stamp Counter}), zwracającą obecną wartość licznika cykli. Licznik służy do precyzyjnego pomiaru czasu i jest inkrementowany co stałą jednostkę czasu, odpowiadającą częstotliwości procesora. Dynamiczne taktowanie obecne we współczesnych rdzeniach nie ma wpływu na częstotliwość tego licznika. W przypadku sprzętu pomiarowego jest to 3.6GHz.
	
	\subsubsection{Implementacja pomiaru}
	W momencie uruchomienia benchmarku ustawiana jest flaga (pkt. 6.4), zapewniająca że komponent wielomianu otrzyma zmianę w każdej klatce, co wywoła pożądaną funkcję. Flaga będzie ustawiana, dopóki benchmark nie zbierze 10000 pomiarów.
	
	\paragraph{Rozpoczęcie pomiaru}
	\begin{verbatim}
		void clock_start() {
		#ifndef __arm64
				start_cycles = __rdtsc();
		#endif
				start = clock();
		}
	\end{verbatim}
	Funkcja ustawia zmienne globalne licznika cykli używając funkcji wewnętrznej \verb|__rdtsc| oraz drugiego licznika używając kanonicznego \verb|clock|. Użycie \verb|__rdtsc| jest wyłączone na platformie ARM, gdzie instrukcja \verb|RDTSC| nie jest dostępna.\\
	

	\paragraph{Zakończenie pomiaru}
	\begin{verbatim}
		void clock_end() {
		#ifndef __arm64
			end_cycles = __rdtsc();
		#endif
			end = clock();
			const clock_t diff = end - start;
			const unsigned long long diff_cycles = end_cycles - start_cycles;
			
			if (samples == -1) return;
			if (samples < CLOCK_SAMPLE_COUNT) {
				samples++;
				sum += diff;
				sum_cycles += diff_cycles;
				if (best == -1 || diff < best) best = diff;
				if (best_cycles == -1 || diff_cycles < best_cycles) best_cycles = diff_cycles;
				
				if (samples == CLOCK_SAMPLE_COUNT) {
					/* wyświetlanie wyniku... */
					samples = -1;
					sum = 0;
					sum_cycles = 0;
					best = -1;
					best_cycles = -1;
				}
			}
		}
	\end{verbatim}
	Liczona jest różnica w z obecnym a poprzednim stanem licznika. Wartość -1 dla \verb|samples| oznacza wyłącony benchmark, a dla \verb|best| i \verb|best_cycles| - brak wartości.
	
	\subsubsection{Narzut pomiaru}
	Aby móc wyciagnąć wnioski z pomiarów sekcji kodu, należy najpierw zmierzyć czas wykonania samego kodu pomiarowego. Wszelkie wyniki zbliżone do narzutu powiaru są niemiarodajne.
	
	\begin{center}
		\begin{tabular}{|c|c|c|c|c|c|} \hline
			\multicolumn{6}{|c|}{Debug} \\\hline
			Pomiar  & Liczba próbek & Średnia RDTSC & Min RDTSC & Średnia clock [ms] & Min clock [ms] \\\hline
			\hline
			1 & 10000 & 3642 & 2592 & 0.000814 & 0.000000 \\\hline
			2 & 10000 & 3561 & 2484 & 0.000791 & 0.000000 \\\hline
			3 & 10000 & 3604 & 2520 & 0.000814 & 0.000000 \\\hline
			5 & 10000 & 3683 & 2592 & 0.000810 & 0.000000 \\\hline
			
		\end{tabular}
	\end{center}
	\begin{center}
		\begin{tabular}{|c|c|c|c|c|c|} \hline
			\multicolumn{6}{|c|}{Release} \\\hline
			Pomiar  & Liczba próbek & Średnia RDTSC & Min RDTSC & Średnia clock [ms] & Min clock [ms] \\\hline
			\hline
			1 & 10000 & 3880 & 2628 & 0.000817 & 0.000000 \\\hline
			2 & 10000 & 3615 & 2556 & 0.000789 & 0.000000 \\\hline
			3 & 10000 & 3713 & 2736 & 0.000796 & 0.000000 \\\hline
			4 & 10000 & 3569 & 2700 & 0.000775 & 0.000000 \\\hline
			
		\end{tabular}
	\end{center}
	Czas jest na tyle mały, że \verb|clock| wskazuje 0 - jest za mało precyzyjny.
	Co ciekawe, wygląda na to, sam narzut w konfiguracji release jest nieco wolniejszy niż w debug.
	
	\pagebreak
	
	\subsection{Pierwotny algorytm}
	\subsubsection{Kod}
	\begin{verbatim}
		void regression(CurvePolynomial* curve, const Vector2* points, const int point_count) {
			clock_start();
			// przygotowanie
			const int degree = point_count - 1;
			const int n = point_count;
			const int m = degree;
			
			float* matrix = malloc((m*m + 2*m) * sizeof (float)  +  m * sizeof (int));
			float* errors = &matrix[m*m];
			float* coeffs = &errors[m];
			int* swap_table = (int*) &coeffs[m];
			
			// wypełananie macierzy i wektora
			for (int row = 0; row < m; ++row) {
				swap_table[row] = row;
				
				for (int col = 0; col < m; ++col) {
					float cell = 0;
					const int power = row + col;
					for (int i = 0; i < n; ++i) {
						cell += pownf(points[i].x, power);
					}
					_get_mat_cell(row, col) = cell;
				}
				float error = 0;
				for (int i = 0; i < n; ++i) {
					error += pownf(points[i].x, row) * points[i].y;
				}
				_get_error(row) = error;
			}
			
			// eliminacja Gaussa - postępowanie proste
			for (int i = 0; i < m - 1; ++i) {
				float max = fabsf(_get_mat_cell(i, i));
				int max_row = i;
				for (int row = i + 1; row < m; ++row) {
					const float current_abs = fabsf(_get_mat_cell(row, i));
					if (current_abs > max) {
						max = current_abs;
						max_row = row;
					}
				}
				if (max_row != i) swap_rows(swap_table, i, max_row);
				
				const float cell = _get_mat_cell(i, i);
				for (int row = i + 1; row < m; ++row) {
					const float coeff = _get_mat_cell(row, i) / cell;
					_get_error(row) -= coeff * _get_error(i);
					for (int col = 0; col < m; ++col) {
						_get_mat_cell(row, col) -= coeff * _get_mat_cell(i, col);
					}
				}
			}
			
			
			// eliminacja Gaussa - postępowanie odweotne
			_get_coeff(m - 1) = _get_error(m - 1) / _get_mat_cell(m - 1, m - 1);
			for (int i = m - 2; i >= 0; --i) {
				float sigma = 0;
				for (int k = i + 1; k < m; ++k) {
					sigma += _get_mat_cell(i, k) * _get_coeff(k);
				}
				_get_coeff(i) = (_get_error(i) - sigma) / _get_mat_cell(i, i);
			}
			
			// przepisanie wyników i zwolnienie pamięci
			for (int i = 0; i < m; ++i) {
				curve->coefficients[i] = _get_coeff(i);
			}
			free(matrix);
			
			curve->order = m;
			curve->start_x = points[0].x;
			curve->end_x = points[point_count - 1].x;
			clock_end();
		}
	\end{verbatim}
	
	\subsubsection{Benchmark}
	\begin{center}
		\begin{tabular}{|c|c|c|c|c|c|} \hline
			\multicolumn{6}{|c|}{Debug} \\\hline
			Pomiar  & Liczba próbek & Średnia RDTSC & Min RDTSC & Średnia clock [ms] & Min clock [ms] \\\hline
			\hline
			1 & 10000 & 1228592 & 1182240 & 0.340353 & 0.328000 \\\hline
			2 & 10000 & 1222491 & 1187676 & 0.338560 & 0.329000 \\\hline
			3 & 10000 & 1221949 & 1184508 & 0.338329 & 0.328000 \\\hline
			4 & 10000 & 1219704 & 1187136 & 0.337932 & 0.328000 \\\hline
			
		\end{tabular}
	\end{center}
	\begin{center}
		\begin{tabular}{|c|c|c|c|c|c|} \hline
			\multicolumn{6}{|c|}{Release} \\\hline
			Pomiar  & Liczba próbek & Średnia RDTSC & Min RDTSC & Średnia clock [ms] & Min clock [ms] \\\hline
			\hline
			1 & 10000 & 951561 & 913932 & 0.263545 & 0.253000 \\\hline
			2 & 10000 & 952509 & 913644 & 0.263831 & 0.253000 \\\hline
			3 & 10000 & 939618 & 913608 & 0.260251 & 0.252000 \\\hline
			4 & 10000 & 942651 & 914544 & 0.261089 & 0.253000 \\\hline
			
		\end{tabular}
	\end{center}
	
	\subsubsection{Wnioski}
	Widać zauważalną różnicę między konfiguracją debug, a release. Wykonanie mieści się poniżej 1 milisekundy, jednak zajmuje dużą ilość czasu w porównaniu do reszty programu oraz powoduje zauważalne spadki liczby klatek na sekundę (z ok. 8000 do ok. 1000).
	
	
\end{document}